%(BEGIN_QUESTION)
% Copyright 2010, Tony R. Kuphaldt, released under the Creative Commons Attribution License (v 1.0)
% This means you may do almost anything with this work of mine, so long as you give me proper credit

Les og lag disposisjon av delen ``Different PID Equations'' i kapitlet ``Closed-Loop Control'' i læreboken din {\it Lessons In Industrial Instrumentation}. Noter sidetallene der viktige illustrasjoner, fotografier, ligninger, tabeller og andre relevante detaljer finnes. Forbered deg på å diskutere gjennomtenkt med lærer og medstudenter de konseptene og eksemplene som behandles i denne lesingen.

\underbar{file i01642}
%(END_QUESTION)





%(BEGIN_ANSWER)


%(END_ANSWER)





%(BEGIN_NOTES)

I parallell-ligningen har hver innstillingskonstant en uavhengig effekt på sin virkning:

$$m = K_p e + {1 \over \tau_i} \int e \> dt + \tau_d {de \over dt} + b \hbox{\hskip 50pt {\bf Parallel PID equation}}$$

\vskip 10pt

I ideal-ligningen påvirker forsterkningskonstanten ($K_p$) alle tre virkningene:

$$m = K_p \left( e + {1 \over \tau_i} \int e \> dt + \tau_d {de \over dt} \right) + b \hbox{\hskip 50pt {\bf Ideal} or {\bf ISA PID equation}}$$

\vskip 10pt

I serie-ligningen påvirker forsterkningskonstanten ($K_p$) alle tre virkningene, mens konstantene $\tau_i$ og $\tau_d$ også påvirker proporsjonalvirkningen:

$$m = K_p \left[ \left({\tau_d \over \tau_i} + 1 \right) e + {1 \over \tau_i} \int e \> dt + \tau_d {de \over dt} \right] + b \hbox{\hskip 25pt {\bf Series} or {\bf Interacting PID equation}}$$

Serieligningen er et resultat av konstruksjon av pneumatiske og analoge elektroniske regulatorer, der det var billigere å bygge en regulator som implementerte denne ligningen enn å implementere en av de to andre ligningene.






\vskip 20pt \vbox{\hrule \hbox{\strut \vrule{} {\bf Forslag til sokratisk diskusjon} \vrule} \hrule}

\begin{itemize}
\item{} Hvorfor bør vi bry oss om at det finnes ulike versjoner av PID-ligningen? Med andre ord: hvilken praktisk betydning har dette for deg som instrumenttekniker?
\item{} Forklar hvordan man matematisk kan utlede hver av formlene for uavhengig virkning (løse for $\Delta m$, bidraget fra hvert reguleringsledd til én utgang) som vises i læreboken.
\item{} Hvilken/hvilke ligning(er) lar oss lage en regulator med bare integralvirkning ved å sette både $K_p$ og $\tau_d$ til null?
\item{} Hvilken/hvilke ligning(er) endrer aggressiviteten til alle tre ledd (P, I og D) når bare $K_p$ justeres?
\item{} Hvorfor har vi i det hele tatt noe så komplekst som ``serie''- eller ``interagerende'' PID-ligning?
\end{itemize}

%INDEX% Reading assignment: Lessons In Industrial Instrumentation, Closed-Loop Control (different PID equations)

%(END_NOTES)
